% This program is free software. It comes without any warranty, to
%   the extent permitted by applicable law. You can redistribute it
%   and/or modify it under the terms of the Do What The Fuck You Want
%   To Public License, Version 2, as published by Sam Hocevar. See
%   http://www.wtfpl.net/ for more details.


% You can choose from the predefined
% \documentclass[
%     nick=robotique,
%     title=Title,
%     subtitle={Subtitle}
% ]{efrei_report_card}

% Or use a custom color
\documentclass[
    name={Custom Name},
    color=ff0000,
    title=Title,
    subtitle={Subtitle}
]{efrei_report_card}

% \usepackage{hyperref}

\author{Jacques Soghomonyan}

\begin{document}
    \maketitle

    \exercise{1 -- Concepts of encapsulation, inheritance, and polymorphism}{
        \question{Encapsulation}{%
            Encapsulation is the way, in Java, the user can assign an \textbf{access modifier} to limit the visibility of certain attributes or methods
        }

        \vspace{2mm}

        \question{Setter and Getter}{
            \begin{multicols}{2}
                \textbf{Setter}\\
                The setter method is a method used to make a private attribute writable by users. 

                \columnbreak

                \textbf{Getter}\\
                The getter method is a method used to make a private attribute readable by users.
            \end{multicols}
        }

        \vspace{2mm}

        \question{this and super}{
            \begin{multicols}{2}
                \textbf{this}\\
                \emph{this} is an implicit parameter passed when calling a method of an object

                \columnbreak

                \textbf{Getter}\\
                The getter method is a method used to make a private attribute readable by users.
            \end{multicols}
        }
    }



\end{document}